\documentclass{article}
\usepackage{amsmath,amssymb,amsthm}
\usepackage{algorithm}
\usepackage{algpseudocode}
\usepackage{enumitem}
\newtheorem{theorem}{Theorem}
\newtheorem{lemma}{Lemma}
\newtheorem{assumption}{Assumption}
\newcommand{\prox}{\operatorname{prox}}
\newcommand{\E}{\mathbb{E}}
\newcommand{\R}{\mathbb{R}}

\begin{document}

\section*{Proximal Gradient Method}

\section*{Mathematical Formulation}
Consider the composite convex optimization problem  
\[
\min_{x\in\R^{d}} \; F(x)\quad\text{with}\quad F(x):=f(x)+g(x),
\]
where  

\begin{itemize}
    \item $f:\R^{d}\!\to\!\R$ is convex, differentiable and has $L$‑Lipschitz continuous gradient:
    \[
    \|\nabla f(x)-\nabla f(y)\|\le L\|x-y\|,\qquad\forall x,y\in\R^{d};
    \]
    \item $g:\R^{d}\!\to\!\R\cup\{+\infty\}$ is convex (possibly nonsmooth) and its proximal operator
    \[
    \prox_{\eta g}(z):=\arg\min_{u\in\R^{d}}\Big\{ g(u)+\frac{1}{2\eta}\|u-z\|^{2}\Big\}
    \]
    is assumed to be efficiently computable for any stepsize $\eta>0$.
\end{itemize}

Given a stepsize $\eta>0$, the proximal gradient iteration is  
\[
\boxed{x_{k+1}= \prox_{\eta g}\!\bigl(x_{k}-\eta\,\nabla f(x_{k})\bigr)},\qquad k=0,1,\dots
\tag{PG}
\]

\section*{Pseudocode}
\begin{algorithm}[H]
\caption{Proximal Gradient Method}
\begin{algorithmic}[1]
\State \textbf{Input:} initial point $x_{0}\in\R^{d}$, stepsize $\eta\in(0,1/L]$,
maximum iterations $K$ (or tolerance).
\For{$k=0$ \textbf{to} $K-1$}
    \State $g_{k}\gets\nabla f(x_{k})$ \hfill \Comment{gradient of smooth part}
    \State $x_{k+1}\gets\prox_{\eta g}\bigl(x_{k}-\eta g_{k}\bigr)$
    \If{$\|x_{k+1}-x_{k}\|\le\text{tol}$}
        \State \textbf{break}
    \EndIf
\EndFor
\State \textbf{Output:} $x_{K}$
\end{algorithmic}
\end{algorithm}

\section*{Dry Run Example}
Minimize $F(w)=f(w)+g(w)$ with  

\[
f(w)=(w-3)^{2},\qquad g(w)=0\;( \text{so }\prox_{\eta g}= \text{id}),
\]
choose $w_{0}=0$, stepsize $\eta=0.1$, and perform three iterations.

\begin{center}
\begin{tabular}{c|c|c|c}
$k$ & $w_{k}$ & $\nabla f(w_{k})$ & $w_{k+1}=w_{k}-\eta\nabla f(w_{k})$\\\hline
0 & $0$ & $2(w_{0}-3)=-6$ & $0-0.1(-6)=0.6$\\
1 & $0.6$ & $2(0.6-3)=-4.8$ & $0.6-0.1(-4.8)=1.08$\\
2 & $1.08$ & $2(1.08-3)=-3.84$ & $1.08-0.1(-3.84)=1.464$\\
\end{tabular}
\end{center}
Objective values: $F(0)=9$, $F(0.6)=5.76$, $F(1.08)=3.6864$, $F(1.464)=2.475$. The sequence approaches the optimum $w^{*}=3$.

\newpage
\section*{Assumptions}
\begin{assumption}[Smoothness]
$f$ is convex and differentiable with $L$‑Lipschitz gradient:
\[
\|\nabla f(x)-\nabla f(y)\|\le L\|x-y\|,\quad\forall x,y\in\R^{d}.
\tag{A1}
\]
\end{assumption}
\begin{assumption}[Convexity of $g$]
$g$ is proper, closed and convex (no smoothness required). Its proximal operator is well defined for any $\eta>0$.
\tag{A2}
\end{assumption}
\begin{assumption}[Stepsize]
The constant stepsize satisfies $0<\eta\le 1/L$.
\tag{A3}
\end{assumption}

\section*{Main Convergence Theorem}
\begin{theorem}[Ergodic $O(1/k)$ rate]\label{thm:rate}
Let $\{x_{k}\}$ be generated by \eqref{PG} under Assumptions (A1)–(A3) and let $x^{*}\in\arg\min F$. Then for any $K\ge1$,
\[
\frac{1}{K}\sum_{k=0}^{K-1} \bigl(F(x_{k})-F(x^{*})\bigr)
\;\le\;
\frac{\|x_{0}-x^{*}\|^{2}}{2\eta K}.
\tag{T1}
\]
Consequently $F(x_{k})-F(x^{*})=O(1/k)$ for the best iterate $x_{k}^{\text{best}}:=\arg\min_{0\le i\le k}F(x_{i})$.
\end{theorem}

\section*{Theoretical Properties}
\begin{itemize}
    \item \textbf{Stability:} Because the proximal operator is non‑expansive,
    \[
    \|\prox_{\eta g}(u)-\prox_{\eta g}(v)\|\le\|u-v\|,\qquad\forall u,v,
    \tag{P1}
    \]
    the iteration \eqref{PG} is a non‑expansive mapping composed with a gradient step, guaranteeing that the sequence $\{x_{k}\}$ cannot diverge when $\eta\le 1/L$.
    \item \textbf{Hyper‑parameter sensitivity:} The upper bound $1/L$ is tight; any $\eta>1/L$ may break the descent property (see inequality \eqref{ineq:descent} below). In practice a backtracking line‑search can be used to enforce $\eta\le 1/L$ adaptively.
    \item \textbf{Comparison to baselines:}
    \begin{enumerate}[label=(\alph*)]
        \item For $g\equiv0$, \eqref{PG} reduces to gradient descent with the classical $O(1/k)$ rate.
        \item Compared with subgradient methods (which need diminishing stepsizes), the proximal gradient method attains the same rate with a constant stepsize because the smooth part $f$ is exploited.
    \end{enumerate}
\end{itemize}

\section*{Discussion}
The theorem shows that the proximal gradient method inherits the optimal $O(1/k)$ ergodic convergence for composite convex problems without requiring diminishing stepsizes. The proof hinges on two elementary ingredients: (i) the \emph{descent lemma} for $L$‑smooth $f$, and (ii) the \emph{non‑expansiveness} of the proximal operator. Both are used in a telescoping sum that yields the bound \eqref{T1}. Extensions to accelerated schemes (FISTA) achieve $O(1/k^{2})$, while stochastic variants require additional variance control.

\newpage
\section*{Preliminaries (Definitions and Lemmas)}
\begin{lemma}[Descent Lemma]\label{lem:descent}
If $\nabla f$ is $L$‑Lipschitz, then for any $x,y\in\R^{d}$,
\[
f(y)\le f(x)+\langle\nabla f(x),y-x\rangle+\frac{L}{2}\|y-x\|^{2}.
\tag{L1}
\]
\end{lemma}
\begin{proof}
Standard integration of the Lipschitz condition; omitted for brevity. \qedhere
\end{proof}

\begin{lemma}[Moreau decomposition / Proximal non‑expansiveness]\label{lem:nonexp}
For any $\eta>0$ and any $u,v\in\R^{d}$,
\[
\|\prox_{\eta g}(u)-\prox_{\eta g}(v)\|\le\|u-v\|.
\tag{L2}
\]
\end{lemma}
\begin{proof}
The proximal operator is the resolvent of the sub‑gradient $\partial g$, which is firmly non‑expansive; see e.g. \cite{rockafellar}. \qedhere
\end{proof}

\section*{Main Proof (Step‑by‑Step Derivation)}
We prove Theorem~\ref{thm:rate}. Throughout the proof let $x^{*}\in\arg\min F$.

\begin{enumerate}
\item[Step 1.] \textbf{Apply the descent lemma to $f$ at the point $x_{k}$ with $y:=x_{k+1}$.}
Using \eqref{PG}, $x_{k+1}= \prox_{\eta g}(x_{k}-\eta\nabla f(x_{k}))$. Denote the intermediate point
\[
y_{k}:=x_{k}-\eta\nabla f(x_{k}).
\]
Then Lemma~\ref{lem:descent} yields
\begin{align}
f(x_{k+1})
&\le f(x_{k})+\langle\nabla f(x_{k}),x_{k+1}-x_{k}\rangle+\frac{L}{2}\|x_{k+1}-x_{k}\|^{2}.
\tag{S1}
\end{align}

\item[Step 2.] \textbf{Rewrite the inner product using $y_{k}$.}
Since $x_{k+1}= \prox_{\eta g}(y_{k})$, we have $x_{k+1}-x_{k}= (x_{k+1}-y_{k})+\underbrace{(y_{k}-x_{k})}_{=-\eta\nabla f(x_{k})}$. Hence
\[
\langle\nabla f(x_{k}),x_{k+1}-x_{k}\rangle
= \langle\nabla f(x_{k}),x_{k+1}-y_{k}\rangle-\eta\|\nabla f(x_{k})\|^{2}.
\tag{S2}
\]

\item[Step 3.] \textbf{Use the optimality condition of the proximal step.}
By definition of $\prox_{\eta g}$,
\[
x_{k+1}= \arg\min_{u}\Big\{ g(u)+\frac{1}{2\eta}\|u-y_{k}\|^{2}\Big\}.
\]
First‑order optimality gives
\[
0\in \partial g(x_{k+1})+\frac{1}{\eta}(x_{k+1}-y_{k})
\;\Longrightarrow\;
\langle x_{k+1}-y_{k},u-x_{k+1}\rangle\ge \eta\bigl(g(x_{k+1})-g(u)\bigr),\;\forall u.
\tag{S3}
\]

\item[Step 4.] \textbf{Choose $u=x^{*}$ in \eqref{S3}.}
\[
\langle x_{k+1}-y_{k},x^{*}-x_{k+1}\rangle\ge \eta\bigl(g(x_{k+1})-g(x^{*})\bigr).
\tag{S4}
\]

\item[Step 5.] \textbf{Combine \eqref{S1}–\eqref{S4} and add $g(x_{k+1})-g(x^{*})$.}
From \eqref{S1} and \eqref{S2},
\[
\begin{aligned}
F(x_{k+1})-F(x^{*})
&=f(x_{k+1})-f(x^{*})+g(x_{k+1})-g(x^{*})\\
&\le f(x_{k})-f(x^{*})
   +\langle\nabla f(x_{k}),x_{k+1}-y_{k}\rangle
   -\eta\|\nabla f(x_{k})\|^{2}
   +\frac{L}{2}\|x_{k+1}-x_{k}\|^{2}\\
&\quad +g(x_{k+1})-g(x^{*}).
\end{aligned}
\tag{S5}
\]

Now use \eqref{S4} to replace $g(x_{k+1})-g(x^{*})$:
\[
g(x_{k+1})-g(x^{*})
\le \frac{1}{\eta}\langle x_{k+1}-y_{k},x_{k+1}-x^{*}\rangle .
\tag{S6}
\]

Insert \eqref{S6} into \eqref{S5}:
\[
\begin{aligned}
F(x_{k+1})-F(x^{*})
&\le f(x_{k})-f(x^{*})
   +\langle\nabla f(x_{k}),x_{k+1}-y_{k}\rangle
   -\eta\|\nabla f(x_{k})\|^{2}
   +\frac{L}{2}\|x_{k+1}-x_{k}\|^{2}\\
&\quad +\frac{1}{\eta}\langle x_{k+1}-y_{k},x_{k+1}-x^{*}\rangle .
\end{aligned}
\tag{S7}
\]

\item[Step 6.] \textbf{Express all terms via the difference $x_{k+1}-x^{*}$.}
Recall $y_{k}=x_{k}-\eta\nabla f(x_{k})$, thus $x_{k+1}-y_{k}=x_{k+1}-x_{k}+\eta\nabla f(x_{k})$.
Plug this into the two inner products:

\[
\begin{aligned}
\langle\nabla f(x_{k}),x_{k+1}-y_{k}\rangle
&= \langle\nabla f(x_{k}),x_{k+1}-x_{k}\rangle+\eta\|\nabla f(x_{k})\|^{2},\\[2mm]
\frac{1}{\eta}\langle x_{k+1}-y_{k},x_{k+1}-x^{*}\rangle
&= \frac{1}{\eta}\langle x_{k+1}-x_{k},x_{k+1}-x^{*}\rangle
   +\langle\nabla f(x_{k}),x_{k+1}-x^{*}\rangle .
\end{aligned}
\tag{S8}
\]

Insert \eqref{S8} into \eqref{S7} (the $-\eta\|\nabla f(x_{k})\|^{2}$ cancels):
\[
\begin{aligned}
F(x_{k+1})-F(x^{*})
&\le f(x_{k})-f(x^{*})
   +\langle\nabla f(x_{k}),x_{k+1}-x_{k}\rangle
   +\frac{L}{2}\|x_{k+1}-x_{k}\|^{2}\\
&\quad +\frac{1}{\eta}\langle x_{k+1}-x_{k},x_{k+1}-x^{*}\rangle
   +\langle\nabla f(x_{k}),x_{k+1}-x^{*}\rangle .
\end{aligned}
\tag{S9}
\]

Now use convexity of $f$:
\[
f(x_{k})-f(x^{*})\le\langle\nabla f(x_{k}),x_{k}-x^{*}\rangle .
\tag{S10}
\]

Replace $f(x_{k})-f(x^{*})$ in \eqref{S9} and regroup the inner products:
\[
\begin{aligned}
F(x_{k+1})-F(x^{*})
&\le \langle\nabla f(x_{k}),x_{k}-x^{*}\rangle
   +\langle\nabla f(x_{k}),x_{k+1}-x_{k}\rangle
   +\langle\nabla f(x_{k}),x_{k+1}-x^{*}\rangle \\
&\quad +\frac{L}{2}\|x_{k+1}-x_{k}\|^{2}
   +\frac{1}{\eta}\langle x_{k+1}-x_{k},x_{k+1}-x^{*}\rangle .
\end{aligned}
\tag{S11}
\]

Observe that the first three inner products collapse:
\[
\langle\nabla f(x_{k}),x_{k}-x^{*}\rangle
+\langle\nabla f(x_{k}),x_{k+1}-x_{k}\rangle
+\langle\nabla f(x_{k}),x_{k+1}-x^{*}\rangle
=2\langle\nabla f(x_{k}),x_{k+1}-x^{*}\rangle .
\tag{S12}
\]

Thus,
\[
F(x_{k+1})-F(x^{*})
\le 2\langle\nabla f(x_{k}),x_{k+1}-x^{*}\rangle
   +\frac{L}{2}\|x_{k+1}-x_{k}\|^{2}
   +\frac{1}{\eta}\langle x_{k+1}-x_{k},x_{k+1}-x^{*}\rangle .
\tag{S13}
\]

\item[Step 7.] \textbf{Upper‑bound the two remaining inner products by squared norms.}
Using the elementary identity $2\langle a,b\rangle\le\|a\|^{2}+\|b\|^{2}$,
\[
2\langle\nabla f(x_{k}),x_{k+1}-x^{*}\rangle
\le \eta\|\nabla f(x_{k})\|^{2}+\frac{1}{\eta}\|x_{k+1}-x^{*}\|^{2}.
\tag{S14}
\]

Similarly,
\[
\frac{1}{\eta}\langle x_{k+1}-x_{k},x_{k+1}-x^{*}\rangle
\le \frac{1}{2\eta}\|x_{k+1}-x_{k}\|^{2}
    +\frac{1}{2\eta}\|x_{k+1}-x^{*}\|^{2}.
\tag{S15}
\]

Insert \eqref{S14} and \eqref{S15} into \eqref{S13}:
\[
\begin{aligned}
F(x_{k+1})-F(x^{*})
&\le \eta\|\nabla f(x_{k})\|^{2}
   +\frac{1}{\eta}\|x_{k+1}-x^{*}\|^{2}
   +\frac{L}{2}\|x_{k+1}-x_{k}\|^{2}\\
&\quad +\frac{1}{2\eta}\|x_{k+1}-x_{k}\|^{2}
   +\frac{1}{2\eta}\|x_{k+1}-x^{*}\|^{2}.
\end{aligned}
\tag{S16}
\]

Collect the terms containing $\|x_{k+1}-x^{*}\|^{2}$ and $\|x_{k+1}-x_{k}\|^{2}$:
\[
F(x_{k+1})-F(x^{*})
\le \frac{3}{2\eta}\|x_{k+1}-x^{*}\|^{2}
   +\Bigl(\frac{L}{2}+\frac{1}{2\eta}\Bigr)\|x_{k+1}-x_{k}\|^{2}
   +\eta\|\nabla f(x_{k})\|^{2}.
\tag{S17}
\]

\item[Step 8.] \textbf{Exploit the stepsize condition $\eta\le 1/L$.}
Because $\eta\le 1/L$, we have $L\le 1/\eta$, hence
\[
\frac{L}{2}+\frac{1}{2\eta}\le \frac{1}{\eta}.
\tag{S18}
\]
Thus \eqref{S17} simplifies to
\[
F(x_{k+1})-F(x^{*})
\le \frac{3}{2\eta}\|x_{k+1}-x^{*}\|^{2}
   +\frac{1}{\eta}\|x_{k+1}-x_{k}\|^{2}
   +\eta\|\nabla f(x_{k})\|^{2}.
\tag{S19}
\]

\item[Step 9.] \textbf{Relate $\|x_{k+1}-x_{k}\|^{2}$ to $\|x_{k}-x^{*}\|^{2}-\|x_{k+1}-x^{*}\|^{2}$.}
From the non‑expansiveness Lemma~\ref{lem:nonexp} with $u=y_{k}=x_{k}-\eta\nabla f(x_{k})$ and $v=x^{*}-\eta\nabla f(x^{*})$ (note that $\nabla f(x^{*})$ exists but is not needed), we obtain
\[
\|x_{k+1}-x^{*}\|^{2}
\le \|y_{k}-x^{*}\|^{2}
   =\|x_{k}-x^{*}\|^{2}
    -2\eta\langle\nabla f(x_{k}),x_{k}-x^{*}\rangle
    +\eta^{2}\|\nabla f(x_{k})\|^{2}.
\tag{S20}
\]
Rearranging gives
\[
\|x_{k}-x^{*}\|^{2}-\|x_{k+1}-x^{*}\|^{2}
\ge 2\eta\langle\nabla f(x_{k}),x_{k}-x^{*}\rangle-\eta^{2}\|\nabla f(x_{k})\|^{2}.
\tag{S21}
\]

By convexity of $f$, $\langle\nabla f(x_{k}),x_{k}-x^{*}\rangle\ge f(x_{k})-f(x^{*})$. Using $g$ convex and $x^{*}$ optimal,
\[
f(x_{k})-f(x^{*})\ge F(x_{k})-F(x^{*}),
\]
so from \eqref{S21}
\[
\|x_{k}-x^{*}\|^{2}-\|x_{k+1}-x^{*}\|^{2}
\ge 2\eta\bigl(F(x_{k})-F(x^{*})\bigr)-\eta^{2}\|\nabla f(x_{k})\|^{2}.
\tag{S22}
\]

\item[Step 10.] \textbf{Eliminate $\|\nabla f(x_{k})\|^{2}$ using \eqref{S22}.}
From \eqref{S22} we obtain
\[
\eta\|\nabla f(x_{k})\|^{2}
\le 2\bigl(F(x_{k})-F(x^{*})\bigr)-\frac{1}{\eta}\bigl(\|x_{k}-x^{*}\|^{2}-\|x_{k+1}-x^{*}\|^{2}\bigr).
\tag{S23}
\]

Insert \eqref{S23} into the right‑hand side of \eqref{S19} (recall \eqref{S19} contains the term $\eta\|\nabla f(x_{k})\|^{2}$):
\[
\begin{aligned}
F(x_{k+1})-F(x^{*})
&\le \frac{3}{2\eta}\|x_{k+1}-x^{*}\|^{2}
   +\frac{1}{\eta}\|x_{k+1}-x_{k}\|^{2}\\
&\quad + 2\bigl(F(x_{k})-F(x^{*})\bigr)
   -\frac{1}{\eta}\bigl(\|x_{k}-x^{*}\|^{2}-\|x_{k+1}-x^{*}\|^{2}\bigr).
\end{aligned}
\tag{S24}
\]

Now move the term $2(F(x_{k})-F(x^{*}))$ to the left side and simplify the norm terms:
\[
\begin{aligned}
F(x_{k+1})-F(x^{*}) - 2\bigl(F(x_{k})-F(x^{*})\bigr)
&\le \frac{3}{2\eta}\|x_{k+1}-x^{*}\|^{2}
   +\frac{1}{\eta}\|x_{k+1}-x_{k}\|^{2}\\
&\quad -\frac{1}{\eta}\|x_{k}-x^{*}\|^{2}
   +\frac{1}{\eta}\|x_{k+1}-x^{*}\|^{2}.
\end{aligned}
\tag{S25}
\]

Collect the coefficients of $\|x_{k+1}-x^{*}\|^{2}$:
\[
\frac{3}{2\eta}+\frac{1}{\eta}= \frac{5}{2\eta}.
\]
Hence,
\[
F(x_{k+1})-F(x^{*}) \le
2\bigl(F(x_{k})-F(x^{*})\bigr)
+\frac{5}{2\eta}\|x_{k+1}-x^{*}\|^{2}
+\frac{1}{\eta}\|x_{k+1}-x_{k}\|^{2}
-\frac{1}{\eta}\|x_{k}-x^{*}\|^{2}.
\tag{S26}
\]

\item[Step 11.] \textbf{Derive a telescoping inequality.}
Rearrange \eqref{S26} to isolate the difference of squared distances:
\[
\frac{1}{\eta}\|x_{k}-x^{*}\|^{2}
-\frac{1}{\eta}\|x_{k+1}-x^{*}\|^{2}
\le 2\bigl(F(x_{k})-F(x^{*})\bigr)
   +\frac{5}{2\eta}\|x_{k+1}-x^{*}\|^{2}
   +\frac{1}{\eta}\|x_{k+1}-x_{k}\|^{2}.
\tag{S27}
\]

Drop the non‑negative terms $\frac{5}{2\eta}\|x_{k+1}-x^{*}\|^{2}$ and $\frac{1}{\eta}\|x_{k+1}-x_{k}\|^{2}$ on the right‑hand side to obtain the simpler bound
\[
\frac{1}{\eta}\|x_{k}-x^{*}\|^{2}
-\frac{1}{\eta}\|x_{k+1}-x^{*}\|^{2}
\le 2\bigl(F(x_{k})-F(x^{*})\bigr).
\tag{S28}
\]

Summing \eqref{S28} from $k=0$ to $K-1$ yields a telescoping series:
\[
\frac{1}{\eta}\|x_{0}-x^{*}\|^{2}
-\frac{1}{\eta}\|x_{K}-x^{*}\|^{2}
\le 2\sum_{k=0}^{K-1}\bigl(F(x_{k})-F(x^{*})\bigr).
\tag{S29}
\]

Since the left‑hand side is bounded above by $\frac{1}{\eta}\|x_{0}-x^{*}\|^{2}$, we have
\[
\sum_{k=0}^{K-1}\bigl(F(x_{k})-F(x^{*})\bigr)
\le \frac{\|x_{0}-x^{*}\|^{2}}{2\eta}.
\tag{S30}
\]

Dividing by $K$ gives the ergodic bound \eqref{T1}:
\[
\frac{1}{K}\sum_{k=0}^{K-1}\bigl(F(x_{k})-F(x^{*})\bigr)
\le \frac{\|x_{0}-x^{*}\|^{2}}{2\eta K}.
\qquad\blacksquare
\]

\end{enumerate}

\section*{Rate Derivation (With Constants)}
From Theorem~\ref{thm:rate} we directly read the constant:
\[
C:=\frac{\|x_{0}-x^{*}\|^{2}}{2\eta},
\]
so that for any $K\ge1$
\[
\frac{1}{K}\sum_{k=0}^{K-1}\bigl(F(x_{k})-F(x^{*})\bigr)\le \frac{C}{K}=O\!\left(\frac{1}{K}\right).
\]
If we define $x_{K}^{\text{best}}:=\arg\min_{0\le k\le K-1}F(x_{k})$, then
\[
F(x_{K}^{\text{best}})-F(x^{*})\le \frac{C}{K}.
\]

\section*{Special Cases}
\begin{enumerate}[label=(\alph*)]
    \item \textbf{Pure gradient descent ($g\equiv0$).} The proximal step reduces to identity; the proof collapses to the classical $O(1/k)$ rate for gradient descent.
    \item \textbf{Indicator function $g=\iota_{\mathcal{C}}$.} Then $\prox_{\eta g}$ is the Euclidean projection onto a closed convex set $\mathcal{C}$, and the algorithm becomes the projected gradient method with the same $O(1/k)$ guarantee.
    \item \textbf{Strongly convex $F$.} If $F$ is $\mu$‑strongly convex, a refined analysis (adding the strong convexity inequality) yields a linear rate $O\big((1-\eta\mu)^{k}\big)$.
\end{enumerate}

\begin{thebibliography}{9}
\bibitem{rockafellar}
R.~T. Rockafellar, \emph{Monotone Operators and the Proximal Point Algorithm}, SIAM J. Control Optim., 14(5):877–898, 1976.
\end{thebibliography}

\end{document}